\documentclass[12pt]{article}

\usepackage[a4paper,margin=2.5cm]{geometry}
\usepackage{setspace}
\usepackage[hidelinks]{hyperref}
\usepackage{longtable}
\usepackage{array}
\usepackage{booktabs}

\title{Literature Review Draft: Continuum-Tipped Minimally Invasive Surgical Robotics}
\author{Aidan Knipe}
\date{\today}

\begin{document}
\onehalfspacing
\maketitle

\begin{abstract}
This draft literature review is structured for a thesis on a UR5e-based macro-micro robotic system with a flexible continuum tip for minimally invasive surgery (MIS). It consolidates relevant findings from prior work in the \textit{Lachlan Report}, identifies which foundations remain valid, and highlights where the state of the art now requires updated evidence. The review prioritizes four technical threads central to this thesis: (1) tendon-driven continuum manipulator design and motion algorithms, (2) integration with a macro robot (UR5e), (3) collision detection and avoidance for safe operation in constrained anatomy, and (4) rapid end-effector exchange mechanisms. It also includes explicit ``missing area'' headings to direct next-stage searching, reading, and synthesis.
\end{abstract}

\section{Project Context and Problem Significance}
This thesis extends prior macro-micro robotic surgery work by coupling a UR5e macro manipulator with a new flexible continuum tip, then improving movement algorithms inherited from earlier implementations. The applied significance is strong: MIS demands dexterity in constrained spaces, low tissue interaction force, high positional repeatability, and robust safety controls.

From the earlier project, the macro-micro concept and teleoperation feasibility are already established at a proof-of-concept level. The key opportunity now is to move from demonstration toward a more clinically credible engineering argument through stronger modeling, safety layers (especially collision handling), and modular hardware architecture (rapid tool exchange).

\section{Inherited Knowledge from Lachlan Report: What Is Still Useful}
The prior report remains highly relevant in these areas:

\subsection{Continuum and Snake-like Design Foundations}
The prior review appropriately positioned tendon-driven continuum/snake-like manipulators as a practical and dominant pathway for surgical dexterity. The classification of continuum architectures and the design trade-off framing (dexterity, stiffness, manufacturability, miniaturization) remain useful to retain.

\subsection{Macro-Micro Integration Rationale}
The macro-micro argument is still strong: a rigid macro arm contributes workspace and gross positioning, while a flexible micro tip contributes distal dexterity and safer interaction in confined regions. This directly aligns with using the UR5e as the macro stage and a continuum tip as the micro stage.

\subsection{Teleoperation and Controller Pipeline Lessons}
The earlier implementation experience (ROS integration, latency considerations, and vibration observations) provides practical system-level lessons for architecture and control-loop design. Even if this thesis shifts focus from pure teleoperation toward autonomy-assisted control, these integration lessons are still relevant.

\subsection{Kinematic Modeling as a Bottleneck}
The prior work highlighted that continuum-tip performance quality is tightly linked to model quality and calibration quality. That finding should be retained as a central thesis narrative point: movement algorithm improvement must be justified through both model structure and validation methodology.

\section{What Must Be Brought Up to Date}
The earlier literature base appears strongest through roughly 2022 and should now be updated in the following directions.

\subsection{Post-2022 Continuum Manipulator Control}
Update with recent results on:
\begin{itemize}
	\item Data-driven and learning-augmented continuum control (model correction, residual learning).
	\item Real-time model-based control under friction, hysteresis, and tendon coupling.
	\item Multi-objective control formulations balancing tracking error, smoothness, and safety constraints.
\end{itemize}

\subsection{Safety-Critical Surgical Robot Control}
The inherited report discusses precision and teleoperation, but safety layers should now be treated as first-class control objectives:
\begin{itemize}
	\item Constraint-aware motion planning and control barrier function style methods.
	\item Contact-aware control and force-limited behavior near anatomical boundaries.
	\item Fault handling and safe-stop behavior for software and communication faults.
\end{itemize}

\subsection{Collision Detection and Avoidance (Core Gap for This Thesis)}
This is one of the most important expansion areas for your project. The review needs a dedicated synthesis of:
\begin{itemize}
	\item Geometric collision detection for continuum bodies (not just rigid-link approximations).
	\item Environment representations (point cloud, mesh, signed distance field) suitable for intraoperative updates.
	\item Real-time avoidance strategies compatible with UR5e + continuum-tip control bandwidth.
\end{itemize}

\subsection{Rapid End-Effector Change and Modular Interfaces}
The prior conclusion already identifies end-effector replacement as cumbersome. This thesis should elevate this into an explicit engineering research stream:
\begin{itemize}
	\item Mechanical quick-change couplers for sterile or semi-sterile workflows.
	\item Repeatable kinematic/electrical interface registration after each swap.
	\item Calibration and safety checks needed after end-effector exchange.
\end{itemize}

\section{Focused Literature Synthesis for Current Thesis Objectives}

\subsection{Thread A: Continuum Tip Design for MIS on UR5e}
The review should compare tendon-driven continuum options based on criteria directly linked to your implementation:
\begin{itemize}
	\item Outer diameter and achievable curvature.
	\item Tendon routing complexity and friction sensitivity.
	\item Manufacturability (3D printing route, assembly complexity, tolerance sensitivity).
	\item Instrument channel requirements and compatibility with intended tasks.
\end{itemize}

\subsection{Thread B: Movement Algorithm Improvement over Lachlan Baseline}
Define baseline algorithm(s) from inherited work, then review candidate improvements under consistent metrics:
\begin{itemize}
	\item Tracking performance (RMS error, peak error, settling behavior).
	\item Robustness to tendon slack, backlash, and parameter drift.
	\item Computational feasibility for real-time execution.
\end{itemize}
The literature section should explicitly map each candidate algorithm class to likely implementation effort and expected gain.

\subsection{Thread C: Collision Detection and Avoidance in Constrained Anatomy}
Your review should separate the problem into three layers:
\begin{itemize}
	\item \textbf{Perception/model layer}: what representation of anatomy and instruments is available in real time?
	\item \textbf{Detection layer}: how are self-collision, tool-environment collision, and forbidden regions detected?
	\item \textbf{Response layer}: how does the controller react (hard stop, velocity scaling, local re-plan, haptic/visual warning)?
\end{itemize}

\subsection{Thread D: Rapid End-Effector Exchange}
This thread should connect mechanism design with control and safety:
\begin{itemize}
	\item Mechanical interchange design options and locking reliability.
	\item Coupling repeatability and resulting pose uncertainty.
	\item Required post-swap validation routine before re-entry into operation.
\end{itemize}

\section{Explicit Missing Research Areas (Add These as Search Targets)}
The following subsection titles are intentionally phrased as missing-area headings to guide your next reading pass.

\subsection{Missing Area 1: Benchmark Datasets and Standardized Evaluation Protocols for Continuum Surgical Robots}
Need: common benchmarks to compare controller improvements credibly across studies.

\subsection{Missing Area 2: Real-Time Collision Avoidance Methods Validated on Macro-Micro Systems with Flexible Distal Tools}
Need: methods tested on systems structurally similar to UR5e + continuum tip.

\subsection{Missing Area 3: Registration Drift and Recalibration After Quick End-Effector Exchange}
Need: evidence on how repeated tool swaps affect kinematic accuracy and safety margins.

\subsection{Missing Area 4: Safety Cases for Human-in-the-Loop MIS Robots with Shared Control}
Need: practical safety argument structures for semi-autonomous assistance during teleoperation.

\subsection{Missing Area 5: Tendon Wear, Sterilization Effects, and Lifecycle Reliability in Reconfigurable Continuum Tools}
Need: long-horizon reliability evidence to support repeated clinical-style usage.

\subsection{Missing Area 6: Task-Specific Performance Metrics Beyond Tip Tracking Error}
Need: clinically meaningful metrics (procedure time, tissue interaction force proxies, collision near-miss counts).

\section{Rubric-Aligned Coverage Notes (MMAN4951/MMAN9451 Interim Report)}

\subsection{Criterion 1: Literature Review (High Weighting)}
To target Distinction/High Distinction bands, this review must do more than summarize papers. It should:
\begin{itemize}
	\item Show conceptual links between manipulator design, control, safety, and modularity.
	\item Explicitly identify and justify knowledge gaps your thesis addresses.
	\item Include recent work (not only foundational papers).
	\item Critically compare methods instead of listing them.
\end{itemize}

\subsection{Criterion 2: Research Question and Plan}
Your final report should ensure the plan is narratively tied to the review:
\begin{itemize}
	\item Each project task should be justified by a gap identified in this section.
	\item Include feasible milestones and contingency options (especially for hardware delays and calibration issues).
\end{itemize}

\subsection{Criterion 3: Project-Dependent Preparations}
The literature should foreshadow practical preparations you can evidence:
\begin{itemize}
	\item UR5e safety/process training,
	\item software stack readiness,
	\item preliminary prototype tests,
	\item risk register for collision and tool-swap failure modes.
\end{itemize}

\subsection{Criterion 4: Document Presentation}
For marks protection, maintain:
\begin{itemize}
	\item consistent citation formatting,
	\item clear figure/table captions,
	\item polished grammar and structured argument flow.
\end{itemize}

\section{Proposed Next Reading and Writing Sequence}
\begin{enumerate}
	\item Build a paper matrix with columns: method class, hardware platform, control rate, collision strategy, modularity strategy, metrics, limitations.
	\item Fill collision-avoidance and quick-change sections first (highest novelty for your topic).
	\item Revisit inherited movement algorithm section and position your planned improvement as a direct response to documented limitations.
	\item Close the review with a concise ``identified gap to thesis objective'' mapping table.
\end{enumerate}

\section{Working Thesis Gap Statement (Draft)}
Current literature and prior project work support feasibility of macro-micro MIS manipulation, but leave a practical translation gap at the intersection of \textbf{accurate continuum movement control}, \textbf{real-time collision-safe operation}, and \textbf{fast, repeatable end-effector reconfiguration}. This thesis addresses that gap on a UR5e-based platform by improving inherited motion algorithms and integrating safety and modularity mechanisms necessary for repeatable experimental deployment.

\end{document}
